\section{A complete space over varying curve backgrounds}
\label{sec:joint-curve-model-spaces}

The polar description also compares models whose backgrounds are different. Keeping the background as a coordinate gives a complete space across the zero-area boundary. The singularity coordinate still records diffuse polar curvature, and taking its atomic part still gives the nearest ideal model.

Fix the compact connected curve $Y$ and a smooth positive form $\eta$ of area one. Let
\[
 \mathcal B(Y)=\{\sigma\in C^\infty(Y,\Lambda^{1,1}_{\mathbb R}):
                         A_\sigma=\textstyle\int_Y\sigma\geq0\}.
\]
Choose $C^k$ norms using a fixed smooth metric, and put
\begin{equation}\label{eq:curve-background-smooth-metric}
 b(\sigma,\tau)=\sum_{k=0}^\infty2^{-k-1}
              \min\{1,\|\sigma-\tau\|_{C^k}\}.
\end{equation}
This is a complete metric for the $C^\infty$ topology on $\mathcal B(Y)$. Indeed, a $b$-Cauchy sequence is Cauchy in every $C^k$ norm. The corresponding limits agree and define a smooth form. Its area is nonnegative by $C^0$ convergence. Conversely, convergence in every norm gives convergence for $b$ by first controlling a finite sum and then its tail.

Let $\mathcal M(\sigma)$ be the ordinary model space of \cref{prop:curve-arbitrary-background-models}, and define
\[
 \mathscr M(Y)=\{(\sigma,p):\sigma\in\mathcal B(Y),
                                  p\in\mathcal M(\sigma)\}.
\]
For such a pair write $\mu_p^\sigma=\mathbf1_{\{p=-\infty\}}(\sigma+\ddc p)$, and set
\begin{equation}\label{eq:joint-curve-model-metric}
 D((\sigma,p),(\tau,q))
       =b(\sigma,\tau)+\|\mu_p^\sigma-\mu_q^\tau\|_{\mathrm{TV}}.
\end{equation}
Let $\mathscr I(Y)$ be its ideal-model locus. Let $\mathscr A_+(Y)$ consist of pairs whose model has analytic singularities and positive non-pluripolar mass in its own background.

\Needspace{20\baselineskip}
\begin{theorem}\label{thm:joint-curve-model-completeness}\index{Model potential!varying curve backgrounds}\index{Metric completeness!joint curve model space}
The metric space $(\mathscr M(Y),D)$ is complete. Its restriction to each fixed-background fiber is the polar metric $d_\sigma$. The map
\[
 R(\sigma,p)=(\sigma,P_\sigma[p]_{\mathcal I})
\]
is a $1$-Lipschitz retraction onto the closed subspace $\mathscr I(Y)$, and gives the unique nearest ideal pair. More precisely,
\begin{equation}\label{eq:joint-curve-analytic-closure}
 \overline{\mathscr A_+(Y)}^{\,D}=\mathscr I(Y),\qquad
 \operatorname{dist}_D((\sigma,p),\mathscr I(Y))
       =(\mu_p^\sigma)^{\mathrm{na}}(Y).
\end{equation}
If $(\sigma_j,p_j)\to(\sigma,p)$ for $D$, then $p_j\to p$ in $L^1(Y,\eta)$ and in capacity for $\eta$, including when $A_\sigma=0$.
\end{theorem}
\begin{proof}
By \cref{prop:curve-arbitrary-background-models}, polar curvature identifies $\mathscr M(Y)$ with the pairs $(\sigma,\mu)$ for which $\sigma\in\mathcal B(Y)$ and $\mu$ is a positive measure carried by a polar set, of mass at most $A_\sigma$. Formula \eqref{eq:joint-curve-model-metric} is the restriction of the product metric $b+\|\cdot\|_{\mathrm{TV}}$, so it separates pairs and restricts to $d_\sigma$ on every fiber.

Suppose $(\sigma_j,\mu_j)$ is Cauchy. Completeness of $b$ gives $\sigma_j\to\sigma$ smoothly, and completeness of finite signed measures gives $\mu_j\to\mu$ in total variation. The limit is positive and satisfies $\mu(Y)\leq A_\sigma$. Choose polar carriers $E_j$ of the $\mu_j$. Their countable union $E$ is polar, and
\[
 \mu(Y\setminus E)\leq\|\mu-\mu_j\|_{\mathrm{TV}}\longrightarrow0.
\]
Thus $\mu$ is polar and gives a model in background $\sigma$. This proves completeness.

Ideal projection takes the atomic part by \cref{prop:curve-arbitrary-background-models}. Atomic extraction contracts total variation, so $R$ is $1$-Lipschitz. The positive atomic measures are closed in total variation: a limit is carried by the countable union of the atom sets of the sequence. Therefore $\mathscr I(Y)$ is closed. Write $\mu=\mu^{\mathrm{at}}+\mu^{\mathrm{na}}$. For every ideal pair $(\tau,q)$ with atomic polar measure $\nu$, mutual singularity gives
\[
 D((\sigma,p),(\tau,q))
 =b(\sigma,\tau)+\mu^{\mathrm{na}}(Y)
                         +\|\mu^{\mathrm{at}}-\nu\|_{\mathrm{TV}}.
\]
Its unique minimum occurs at $\tau=\sigma$ and $\nu=\mu^{\mathrm{at}}$. This proves the retraction and distance assertions.

For $A_\sigma>0$, finite atomic measures with rational coefficients and total mass strictly below $A_\sigma$ approximate every positive atomic measure of mass at most $A_\sigma$ in total variation. Their models in background $\sigma$ are analytic by the curve classification and smooth-gauge transport, and have positive residual mass. If $A_\sigma=0$, the ideal pair has zero polar measure. The pairs with backgrounds $\sigma+j^{-1}\eta$ and zero polar measure converge to it for $D$. These models are bounded, hence have trivial analytic singularities, and their residual masses are $j^{-1}>0$. Every ideal pair is therefore in the closure of $\mathscr A_+(Y)$. Conversely, an analytic model is its ideal projection by \cref{prop:analytic-envelope-type}; closedness of $\mathscr I(Y)$ proves \eqref{eq:joint-curve-analytic-closure}.

Finally let $(\sigma_j,p_j)\to(\sigma,p)$ and write $\mu_j=\mu_{p_j}^{\sigma_j}$ and $\mu=\mu_p^\sigma$. If $A_\sigma>0$, put $\lambda_j=A_{\sigma_j}/A_\sigma$. For all sufficiently large $j$ these numbers are positive and tend to one, and
\[
 \|\lambda_j^{-1}\mu_j-\mu\|_{\mathrm{TV}}
 \leq\lambda_j^{-1}\|\mu_j-\mu\|_{\mathrm{TV}}
                  +|\lambda_j^{-1}-1|\,\mu(Y)\longrightarrow0.
\]
Apply \cref{cor:curve-varying-background-continuity}. If $A_\sigma=0$, the background convergence is uniform and all areas are nonnegative. The uniform collapse of \cref{thm:curve-zero-area-collapse} gives $L^1$ and capacity convergence to the unique normalized potential in background $\sigma$.
\end{proof}

The $C^\infty$ background topology is part of this completeness statement. A Cauchy sequence of smooth forms for the uniform norm can instead have a nonsmooth limit. At a fixed positive-area background, $D$ induces the positive-volume topology of \cref{thm:positive-model-complete-metric}, while allowing sequences to converge to zero-mass models. It therefore does not retain the reciprocal-volume Cauchy condition of that particular complete metric.

The joint metric also makes smooth gauges and positive scaling compatible across all fibers.

\Needspace{17\baselineskip}
\begin{corollary}\label{cor:joint-curve-model-gauge}\index{Model potential!joint gauge isometry}
For $f\in C^\infty(Y,\mathbb R)$, the map
\[
 \mathcal G_f(\sigma,p)
   =(\sigma+\ddc f,P_{\sigma+\ddc f}[p-f])
\]
is an isometry of $\mathscr M(Y)$, with inverse $\mathcal G_{-f}$ and $\mathcal G_g\mathcal G_f=\mathcal G_{f+g}$. It preserves analytic models and commutes with $R$.

For $\lambda>0$, the bijection $\mathcal H_\lambda(\sigma,p)=(\lambda\sigma,\lambda p)$ commutes with $R$ and satisfies
\[
 \min\{1,\lambda\}D(x,y)
 \leq D(\mathcal H_\lambda x,\mathcal H_\lambda y)
 \leq\max\{1,\lambda\}D(x,y).
\]
\end{corollary}
\begin{proof}
Smooth-gauge transport in \cref{prop:curve-arbitrary-background-models} preserves the polar measure and the ordinary type after subtraction of $f$. In polar coordinates, $\mathcal G_f$ is therefore $(\sigma,\mu)\mapsto(\sigma+\ddc f,\mu)$. Its area is unchanged, and $b(\sigma+\ddc f,\tau+\ddc f)=b(\sigma,\tau)$. This proves the isometry, inverse and composition laws. Atomic extraction commutes with this map, and smooth subtraction and bounded changes preserve analytic singularities.

Scaling the ordinary-envelope competitors shows that $\lambda p$ is a normalized model for $\lambda\sigma$. In polar coordinates this is $(\sigma,\mu)\mapsto(\lambda\sigma,\lambda\mu)$, with inverse scaling by $\lambda^{-1}$. Atomic extraction commutes with scaling. For $r\geq0$,
\[
 \min\{1,\lambda\}\min\{1,r\}
 \leq\min\{1,\lambda r\}
 \leq\max\{1,\lambda\}\min\{1,r\}.
\]
Apply this to every $C^k$ norm in \eqref{eq:curve-background-smooth-metric}; the total-variation term is multiplied by $\lambda$. The two bounds follow.
\end{proof}
